\section{Plausibility-proxy rubric and its sensitivity}\label{supp:plausibility}

The plausibility axis of the leaderboard is a documented, tradition-level proxy, not a
human-subjects measurement. This section records how each value was assigned. It also shows
that the headline contrast does not depend on the particular numbers chosen. The Methods
(\Sectionname\ref{sec:triad}) name two reporting axes. Faithfulness is measured on the machine against the
intervention oracle of \Sectionname\ref{sec:oracle}. The plausibility proxy stands in for how convincing a
tradition's explanations are usually taken to be. We assign the proxy once per method
tradition, before any faithfulness score on our machine is computed, so it cannot be tuned to
the result. Its only job is to expose the danger zone of explanations that look convincing yet
are unfaithful. For that job it must be coarse, fixed in advance, and sourced rather than
precise.

The rubric assigns one proxy per tradition on a fixed five-step scale. It scores two
documented properties of each tradition's typical output and takes their average. The first
property is whether the explanation produces a dense, picture-like artifact that a reader sees
as an account of the computation. Heatmaps and saliency-style maps over the input are read as
convincing by construction. This is the perceptual-immediacy axis. The second property is
whether the field has historically treated the tradition's output as a self-evident
explanation rather than a hypothesis to be tested. This is the adoption-as-explanation axis. A
tradition that scores high on both, such as gradient saliency, receives a high proxy. A
tradition whose output is an austere intervention table or a circuit edge list, rarely
mistaken for a finished story, receives a low one. The scale is deliberately blunt,
$\{0.3,0.5,0.55,0.75,0.8,0.85,0.9\}$. A proxy meant to expose a danger zone gains nothing from
false precision, and a finer scale would invite over-reading. Every value in
Table~\ref{supp:tab:plausrubric} traces to the \jkey{plausibility_proxy} field of the
per-method rows in \texttt{compare/out/leaderboard.json}. The table only groups those rows by
tradition and records the rationale.

\begin{table}[t]
\centering
\caption{\textbf{Plausibility-proxy rubric.} The proxy assigned to each method tradition,
with the documented rationale. Values are the \jkey{plausibility_proxy} field of the
leaderboard rows (\texttt{compare/out/leaderboard.json}). They are tradition-level and fixed
before any faithfulness score is computed, so the proxy cannot be tuned to the measured
faithfulness. A higher value means an explanation that the field tends to find more
convincing, not one that is more faithful. The oracle (\Sectionname\ref{sec:oracle}) is the positive control and
is not a method under test.}
\label{supp:tab:plausrubric}
\begin{tabular}{@{}llp{6.4cm}@{}}
\toprule
Tradition & Proxy & Rationale (documented basis) \\
\midrule
gradient        & $0.90$ & Dense input-space heatmaps read as a visual account of the
                            decision; the most widely adopted XAI outputs, routinely shown
                            as if self-explanatory. \\
probing         & $0.85$ & A trained decoder that ``finds'' a concept reads as evidence the
                            concept is used; high adoption, though decodability is not
                            causal use. \\
correlational   & $0.80$ & Tuning curves and coupling spectra match the neuroscience idiom
                            of a convincing single-unit account. \\
dim.\ reduction & $0.75$ & Low-dimensional factors and dictionary atoms look like a clean
                            summary of the representation. \\
intervention    & $0.55$ & Occlusion and counterfactual outputs are quantitative tables, less
                            immediately picture-like, treated as tests rather than stories. \\
causal          & $0.50$ & Circuit edge lists and patching results are austere artifacts
                            rarely mistaken for a finished narrative. \\
descriptive     & $0.30$ & Whole-state recording is an explicit non-explanation: it withholds
                            nothing and reads as a dump, not an account. \\
\midrule
oracle (control)& $1.00$ & The \Sectionname\ref{sec:oracle} ground truth by construction; positive control, not a
                            method under test. \\
\bottomrule
\end{tabular}
\end{table}

The paper draws a contrast between this proxy and measured faithfulness. On the committed
records the two move in opposite directions. Across the thirty methods under test the Spearman
rank correlation between faithfulness and the proxy is $-0.36$, a negative association.
The traditions the field finds most convincing are, on this machine, among the least faithful.
Eight methods fall in the danger zone, defined here as a proxy of at least $0.70$ together
with a faithfulness below $0.40$. They include the correlational neuroscience methods, Expected
Gradients, guided backprop, linear probing, and the sparse autoencoder. The per-method extreme
makes the point sharply. Vanilla saliency has the highest proxy, $0.9$, yet scores faithfulness
$0.115$ on the position regime. Activation patching has a lower proxy, $0.5$, and scores $1.0$.
Here the less faithful method is the more plausible one. These figures are produced by
\texttt{tools/xai\_study/compare/plausibility\_sensitivity.py} from the leaderboard and
recorded in \texttt{compare/out/plausibility\_sensitivity.csv}.

A proxy fixed by a coarse rubric raises the question of whether the contrast is an artifact of
the particular numbers. The answer is that it is not. We perturb the proxy under four
deterministic schemes and recompute the relationship each time. The analysis uses fixed,
index-keyed offsets rather than random draws, so it re-derives bit-for-bit. The first scheme
adds a fixed per-method jitter of magnitude $\sigma\in\{0.05,0.1,0.2\}$, alternating in sign by
row. The second compresses every proxy toward the common mean by a fraction $\sigma$, which
preserves the rank order of traditions. The third uses the same index-keyed offsets at
$\sigma\in\{0.1,0.2\}$, large enough to swap neighbouring traditions' ranks. The fourth drops
one whole tradition at a time and recomputes on the remainder. Across all fourteen
perturbations the rank correlation between faithfulness and proxy stays negative. It ranges
from $-0.76$ to $-0.08$. The danger zone stays populated, with
between four and ten methods. The per-method anchor sign holds in every scheme where both
anchor methods are present, the more plausible saliency remaining the less faithful. Two runs
bound the range. The strongest is the leave-the-gradient-tradition-out run. It removes the most
plausible and least faithful block at once, and still leaves six danger-zone methods and a rank
correlation of $-0.76$. The weakest is the leave-the-causal-tradition-out run, where the
correlation falls to $-0.08$, close to zero; removing the faithful causal block flattens the
association, so the contrast is carried in part by that block. Table~\ref{supp:tab:plaussens}
reports the full sweep. The contrast keeps its sign under every alternative proxy assignment we
tried: plausibility and faithfulness are not the same axis, and high plausibility is no
guarantee of faithfulness.

\begin{table}[t]
\centering
\caption{\textbf{Proxy-sensitivity sweep.} The faithfulness-versus-proxy relationship under
each perturbation, from \texttt{compare/out/plausibility\_sensitivity.csv}. ``$\rho$'' is the
Spearman rank correlation between faithfulness and proxy over the methods present. ``danger''
is the count of high-proxy ($\geq0.70$), low-faithfulness ($<0.40$) methods. ``anchor'' is
$1$ when the more plausible saliency is the less faithful of the saliency/activation-patching
pair, and blank where a leave-one-out run removes an anchor. The baseline rubric is the first
row. Across every scheme the correlation stays negative and the danger zone stays populated.}
\label{supp:tab:plaussens}
\begin{tabular}{@{}llrrrc@{}}
\toprule
Scheme & Param & $n$ & $\rho$ & Danger & Anchor \\
\midrule
baseline                & ---           & 30 & $-0.36$ &  8 & 1 \\
additive jitter         & $\sigma=0.05$ & 30 & $-0.40$ &  8 & 1 \\
additive jitter         & $\sigma=0.10$ & 30 & $-0.41$ &  8 & 1 \\
additive jitter         & $\sigma=0.20$ & 30 & $-0.42$ & 10 & 1 \\
rank-preserving         & $\sigma=0.10$ & 30 & $-0.36$ &  8 & 1 \\
rank-preserving         & $\sigma=0.20$ & 30 & $-0.36$ &  8 & 1 \\
rank-jittering          & $\sigma=0.10$ & 30 & $-0.41$ &  8 & 1 \\
rank-jittering          & $\sigma=0.20$ & 30 & $-0.42$ & 10 & 1 \\
leave-out causal        & ---           & 24 & $-0.08$ &  8 & \\
leave-out correlational & ---           & 26 & $-0.40$ &  4 & 1 \\
leave-out descriptive   & ---           & 29 & $-0.29$ &  8 & 1 \\
leave-out dim.\ red.    & ---           & 27 & $-0.34$ &  7 & 1 \\
leave-out gradient      & ---           & 20 & $-0.76$ &  6 & \\
leave-out intervention  & ---           & 25 & $-0.34$ &  8 & 1 \\
leave-out probing       & ---           & 29 & $-0.36$ &  7 & 1 \\
\bottomrule
\end{tabular}
\end{table}
